\documentclass[11pt]{article}
\usepackage[a4paper,margin=1in]{geometry}
\usepackage{fontspec}
\usepackage[boldfont=false]{xeCJK}
\setCJKmainfont{FandolSong-Regular.otf}
\usepackage{amsmath}
\usepackage{amssymb}
\usepackage{array}
\usepackage{booktabs}
\usepackage{graphicx}
\usepackage{adjustbox}
\usepackage{enumitem}
\setlist[itemize]{topsep=2pt,partopsep=0pt,parsep=0pt,itemsep=2pt,leftmargin=1.4em}
\setlist[enumerate]{topsep=2pt,partopsep=0pt,parsep=0pt,itemsep=2pt,leftmargin=1.6em}
\usepackage{parskip}
\usepackage{fvextra}
\fvset{breaklines=true,breakanywhere=true,fontsize=\small}
\setlength{\parindent}{0pt}

\begin{document}

\begin{center}
  {\LARGE\bfseries Government microeconomic intervention}\\[0.35em]
  {\large A Level Economics}
\end{center}
\vspace{0.6em}

\section*{Why markets fail}

A free market often gives a good result, but not always. \textbf{market failure} 市场失灵 is when the market does not \textbf{allocate} 配置 (share out) resources efficiently — too much or too little of a good is made. Then the government may step in.

There are several reasons for market failure.

\subsection*{Externalities}

An \textbf{externality} 外部性 is a cost or a benefit that falls on a \textbf{third party} 第三方 — someone who is not the buyer or the seller.

\begin{itemize}
\item a \textbf{negative externality} 负外部性 is an \textbf{external cost} 外部成本 paid by others, like the pollution from a factory. The \textbf{social cost} 社会成本 (the cost to all of society) is bigger than the \textbf{private cost} 私人成本 (the cost to the producer). When people ignore the external cost, the market makes \textbf{too much} of the good.

\item a \textbf{positive externality} 正外部性 is an \textbf{external benefit} 外部收益 enjoyed by others, like the protection that vaccination gives to everyone. Here the market makes \textbf{too little}, because buyers ignore the benefit to others.

\end{itemize}

\subsection*{Public goods}

A \textbf{public good} 公共物品 (like street lighting or national defence) has \textbf{non-rivalry} 非竞争性 and \textbf{non-excludability} 非排他性. Because non-payers cannot be shut out, there is a \textbf{free-rider problem} 搭便车问题, so private firms make no profit and supply none. The market fails completely.

\subsection*{Merit and demerit goods}

A \textbf{merit good} 优值品 (like education) is under-consumed; a \textbf{demerit good} 劣值品 (like cigarettes) is over-consumed. People misjudge them because of \textbf{information failure} 信息失灵 — they do not know, or do not use, the true benefit or harm.

\subsection*{Monopoly power}

\textbf{monopoly} 垄断 means one firm dominates a market. A firm with \textbf{monopoly power} 垄断势力 can cut output and raise the price above the competitive level. This wastes resources and harms consumers.

\subsection*{Inequality and asymmetric information}

\begin{itemize}
\item the market rewards people who own resources, so it can produce great \textbf{inequality} 不平等 — some people are left with too little to live on.

\item \textbf{asymmetric information} 信息不对称 is when one side of a deal knows more than the other (a used-car seller knows the faults; the buyer does not). This leads to poor decisions and unfair deals.

\end{itemize}

\section*{Methods of government intervention}

\subsection*{Indirect taxes}

An \textbf{indirect tax} 间接税 is a tax on spending. It raises a firm's costs, shifts the supply curve to the left, and so raises the price and cuts the quantity. Governments tax demerit goods (like tobacco) and goods with negative externalities to make people use less.

The \textbf{incidence of a tax} 税收归宿 is who really pays it. This depends on elasticity:

\begin{itemize}
\item if demand is \textbf{inelastic} 缺乏弹性, consumers pay most of the tax (the price rises by nearly the full tax).

\item if demand is \textbf{elastic} 富有弹性, producers pay most of the tax (they cannot pass it on).

\end{itemize}

\subsection*{Subsidies}

A \textbf{subsidy} 补贴 is a government payment to producers. It lowers their costs, shifts the supply curve to the right, and so lowers the price and raises the quantity. Subsidies are used for merit goods and goods with positive externalities.

\subsection*{Maximum and minimum prices}

\begin{itemize}
\item a \textbf{maximum price} 最高限价 is a legal top price, set \textbf{below} equilibrium (for example, rent control). Because the price is held down, quantity demanded is greater than quantity supplied, so there is a \textbf{shortage} 短缺.

\item a \textbf{minimum price} 最低限价 is a legal bottom price, set \textbf{above} equilibrium (for example, a minimum wage, or a floor for farm prices). Because the price is held up, quantity supplied is greater than quantity demanded, so there is \textbf{excess supply} 超额供给.

\end{itemize}

\subsection*{Buffer stocks}

A \textbf{buffer stock} 缓冲库存 is used to steady the price of farm goods. The government buys and stores the good when the price is low, and sells from the store when the price is high. This keeps the price inside a target band.

\subsection*{Regulation, provision and permits}

\begin{itemize}
\item \textbf{regulation} 管制 means rules and laws — bans, age limits, safety standards, or pollution limits.

\item \textbf{state provision} 政府提供 means the \textbf{government} 政府 supplies a good directly, often free at the point of use (public goods and merit goods like schools and hospitals).

\item \textbf{tradable permits} 可交易许可证 set a total cap on pollution and give firms permits to pollute. Firms that pollute less can sell spare permits to firms that pollute more. This caps total pollution at the lowest cost.

\end{itemize}

\subsection*{Government failure}

Intervention does not always help. \textbf{government failure} 政府失灵 is when the government's action makes the use of resources \textbf{worse}, not better. Common causes:

\begin{itemize}
\item poor information, so the government sets the wrong tax or subsidy.

\item high cost of running the policy.

\item unintended results — for example, a maximum price can create an illegal \textbf{black market} 黑市.

\item decisions made for political reasons rather than for efficiency.

\end{itemize}

\section*{Income and wealth inequality}

It helps to separate two ideas:

\begin{itemize}
\item \textbf{income} 收入 is a \textbf{flow} 流量 — money received over a period, such as wages, rent and interest.

\item \textbf{wealth} 财富 is a \textbf{stock} 存量 — the value of assets owned at a point in time, such as houses, shares and savings.

\end{itemize}

Wealth produces income (savings earn interest), and income can be saved to build wealth, so the two are linked. Inequality has many causes: differences in skills and wages, the ownership of wealth, inheritance, and unemployment.

\subsection*{Policies to reduce inequality}

\begin{itemize}
\item a \textbf{progressive tax} 累进税 takes a larger \textbf{percentage} 百分比 from higher incomes than from lower incomes. This narrows the gap after tax.

\item \textbf{benefits} 福利金 (also called \textbf{transfer payments} 转移支付) are money paid by the government to the poor, the old, and the unemployed.

\item \textbf{state provision} of free services, like education and health care, raises the real living standards of poorer people.

\end{itemize}

Together these policies bring about \textbf{redistribution} 再分配 — moving income and wealth from richer to poorer people. But there is a \textbf{trade-off} 取舍: very high taxes and large benefits may weaken the \textbf{incentive} 激励 to work and to invest.


\end{document}
